\documentclass[border=10pt]{standalone}
\usepackage{fontspec}
\AtBeginDocument{\hyphenation{zenodo}}% never hyphenate "zenodo" inside a DOI
\setmainfont{Noto Serif}[Ligatures=TeXOff]
\XeTeXgenerateactualtext=1
\newfontfamily\bengalifont{Noto Serif Bengali}[Script=Bengali, NFSSFamily=bnfam]
\usepackage[Latin,Bengali]{ucharclasses}
% group-free font switching (as in header-common.tex): robust inside braces, cells and quotes
\makeatletter
\def\BN@saved{\rmdefault}
\def\BN@bn{bnfam}
\newcommand\BNin{\edef\BN@tmp{\f@family}\ifx\BN@tmp\BN@bn\else\let\BN@saved\BN@tmp\fi\fontfamily{bnfam}\selectfont}
\newcommand\BNout{\fontfamily{\BN@saved}\selectfont}
\makeatother
\setTransitionsFor{Bengali}{\BNin}{\BNout}
% Typographic quotes and dashes are not in a ucharclasses class (class 0), so a Bangla word right
% after “ or ‘ would not trigger the Bangla font. Treat them as Latin so that every switch fires.
\XeTeXcharclass"2018=\XeTeXcharclass`A
\XeTeXcharclass"2019=\XeTeXcharclass`A
\XeTeXcharclass"201C=\XeTeXcharclass`A
\XeTeXcharclass"201D=\XeTeXcharclass`A
\XeTeXcharclass"2013=\XeTeXcharclass`A
\XeTeXcharclass"2014=\XeTeXcharclass`A
\XeTeXcharclass"2026=\XeTeXcharclass`A
\usepackage{newunicodechar}
\newunicodechar{।}{{\bengalifont\char"0964}}
\newunicodechar{॥}{{\bengalifont\char"0965}}

\usepackage[table]{xcolor}
\usepackage{booktabs,tabularx,array}
\definecolor{bapilrule}{HTML}{1F2937}
\definecolor{bapilstripe}{HTML}{F3F4F6}
\renewcommand{\arraystretch}{1.28}
\setlength{\tabcolsep}{5pt}

\begin{document}
\begin{minipage}{17cm}
{\large\bfseries সারণি ২৬।}\ {\large নির্দেশক সংরক্ষণ-সূচি}\par\vspace{7pt}
\small
\rowcolors{2}{bapilstripe}{white}
\begin{tabularx}{\linewidth}{@{}>{\hsize=1.040\hsize\raggedright\arraybackslash}X>{\hsize=0.920\hsize\raggedright\arraybackslash}X>{\hsize=1.040\hsize\raggedright\arraybackslash}X@{}}
\arrayrulecolor{bapilrule}\toprule
\rowcolor{white}\textbf{রেকর্ডের ধরন} & \textbf{সংরক্ষণকাল} & \textbf{নিষ্পত্তি বা সংরক্ষণাগার} \\
\midrule
মূল উন্মুক্ত নথি (DC0) & অনির্দিষ্টকাল & চেকসামসহ সংরক্ষিত; বাতিল সংস্করণও রাখা হয় \\
নিষ্কাশিত টেক্সট ও সূচি & প্রতিস্থাপিত না হওয়া পর্যন্ত & মূল নথি থেকে পুনর্নির্মাণ; পুনরুৎপাদনের জন্য পুরোনো সূচি এক বছর রাখা হয় \\
নিরীক্ষা লগ (প্রশ্ন, উৎস, ফলাফল) & ৫ বছর & শুধু-পঠন আকারে সংরক্ষণ; এরপর মুছে ফেলা \\
লগে প্রশ্নের টেক্সট & ২ বছর & সামষ্টিক পরিসংখ্যানে সংক্ষেপিত \\
মূল পরামর্শ-মতামত (DC2) & পরিচয়-অপসারণের ৯০ দিনের মধ্যে মুছে ফেলা & সনদসহ নিরাপদ অপসারণ \\
পরিচয়-অপসারিত বিশ্লেষণ-ডেটা & পাইলটের মেয়াদ এবং আরও ২ বছর & চুক্তি অনুযায়ী অপসারণ বা সংরক্ষণ \\
মডেল উপকরণ ও মূল্যায়ন প্রতিবেদন & ব্যবস্থার আয়ুষ্কাল এবং আরও ৩ বছর & লাইসেন্স অনুমতি দিলে Zenodo-তে সংরক্ষণ \\
ঘটনার রেকর্ড ও ERRC-এর সিদ্ধান্ত & ৭ বছর & সারাংশ প্রকাশিত; পূর্ণ রেকর্ড সংরক্ষিত \\
\bottomrule
\end{tabularx}
\par\vspace{5pt}{\footnotesize\itshape ব্যক্তিগত উপাত্ত সুরক্ষা আইন, ২০২৬, জাতীয় উপাত্ত ব্যবস্থাপনা আইন, ২০২৬ ও এগুলোর অধীনে জারিকৃত বিধিমালার আলোকে আইনি পর্যালোচনা সাপেক্ষ।}
\par\vspace{3pt}{\scriptsize\color{gray} উৎস: BAPIL টেকনিক্যাল ব্লুপ্রিন্ট, সংস্করণ ১.০ (২০২৬)। CC BY 4.0।}
\end{minipage}
\end{document}
